% Hodina 66, rozpis řádek 66. Zdroje: učebnice
% ucebnice/02_Funkce_a_rovnice/Linearni_funkce_a_rovnice/Lineární funkce.tex,
% části Lineární funkce a konstantní případ; lokální Realisticky.cz
% 020111_Linearni_funkce_IV/020111_Lekce.pdf. Příklady a graf vlastní.
\documentclass[11pt,a4paper]{article}
% Spolecny styl pro ucitelske pripravy MAT (Work/MAT)
% Verze 1.1, 2026-10-23
\NeedsTeXFormat{LaTeX2e}
\ProvidesPackage{priprava}[2026/10/23 v1.1 Sablona ucitelskych priprav MAT]

% Zakladni typografie a cestina
\RequirePackage[T1]{fontenc}
\RequirePackage[utf8]{inputenc}
\RequirePackage[czech]{babel}
\RequirePackage{lmodern}
\RequirePackage{microtype}
\RequirePackage{amsmath,amssymb}
\RequirePackage{xcolor}
\RequirePackage{enumitem}
\RequirePackage[most]{tcolorbox}
\RequirePackage{tabularx,booktabs,array}
\RequirePackage{fancyhdr}
\RequirePackage{needspace}
\RequirePackage{geometry}
\RequirePackage[hidelinks]{hyperref}

\geometry{a4paper,top=17mm,bottom=18mm,left=23mm,right=23mm}
\definecolor{notegray}{gray}{0.39}
\setlength{\parindent}{0pt}
\setlength{\parskip}{2.5pt plus 1pt minus 0.5pt}
\setlist[itemize]{topsep=2pt,itemsep=1pt,parsep=0pt,leftmargin=1.45em}
\setlist[enumerate]{topsep=2pt,itemsep=2pt,parsep=0pt,leftmargin=1.65em}
\setcounter{secnumdepth}{0}
\renewcommand{\arraystretch}{1.12}

% Zahlavi: vlevo rocnik | globalni tema, vpravo cislo: nazev hodiny
% Zapati zustava zcela prazdne.
\pagestyle{fancy}
\fancyhf{}
\renewcommand{\headrulewidth}{0.35pt}
\renewcommand{\footrulewidth}{0pt}
\setlength{\headheight}{14pt}
\newcommand{\nastavPripravu}[4]{%
  \fancyhead[L]{\small #1\enspace\textbar\enspace #2}%
  \fancyhead[R]{\small #3: #4}%
}

% Minimalisticke cernobile boxy. Kazda definice zvlast.
\tcbset{commonlesson/.style={enhanced,breakable=false,colback=white,colframe=black,
  colbacktitle=white,coltitle=black,boxrule=0.45pt,
  arc=0.8mm,outer arc=0.8mm,left=3mm,right=3mm,top=1.6mm,bottom=1.6mm,
  toptitle=1mm,bottomtitle=0.7mm,fonttitle=\bfseries,
  attach boxed title to top left={xshift=2.5mm,yshift=-1.15mm},
  boxed title style={colback=white,colframe=white,boxrule=0pt,left=1mm,right=1mm}}}
\newtcolorbox{zadaniBox}[1][]{commonlesson,title={Zadání},#1}
\newtcolorbox{prikladBox}[1][]{commonlesson,title={Příklad},#1}
\newtcolorbox{definiceBox}[1][]{commonlesson,title={Definice},#1}

% Spolecny box pro zadani a jeho reseni, bez mezery mezi castmi.
% Pouziti:
%   \begin{prikladReseniBox}
%     Zadani ulohy.
%     \tcblower
%     \begin{align*} ... \end{align*}
%   \end{prikladReseniBox}
% Spodni cast ma jemny sedy podklad a delici linku; lze v ni sazet
% vice radku matematiky. Obe casti maji jen jeden spolecny ramecek.
\tcbset{solutionlesson/.style={
  commonlesson,
  enhanced jigsaw,
  breakable,
  colbacklower=black!3,
  segmentation style={solid,draw=black!25,line width=0.35pt},
  fontlower=\small\color{notegray},
  before lower={\noindent\textbf{Řešení:}\enspace},
  leftlower=3mm,rightlower=3mm
}}
\newtcolorbox{prikladReseniBox}[1][]{solutionlesson,title={Příklad},#1}
\newtcolorbox{zadaniReseniBox}[1][]{solutionlesson,title={Zadání},#1}

% Pomucky pro autora pripravy: sedive poznamky; \reseni pro starsi pripravy
\newcommand{\poznamka}[1]{\par\begingroup\small\color{notegray}\itshape
  \textnormal{\textbf{Poznámka pro učitele.}}\enspace #1\par\endgroup}
\newcommand{\reseni}[1]{\par\begingroup\small\color{notegray}
  \textbf{Očekávané řešení:}\enspace #1\par\endgroup}
\newcommand{\blok}[1]{\needspace{5\baselineskip}\subsection*{#1}}
\newcommand{\vyrok}[1]{\uv{#1}}

\usepackage{tikz}
\nastavPripravu{1. ročník}{Funkce}{66}{Konstantní funkce}
\begin{document}
\blok{Návaznost}
Při konstrukci $f(x)=ax+b$ jsme v hodině 65 používali svislou změnu $\Delta y=a\Delta x$. Co se stane při $a=0$? Každá změna argumentu dává $\Delta y=0$. Výstup se nemění, i když se vstup mění.
\begin{definiceBox}
\textbf{Konstantní funkce} na oboru $D$ má předpis $f(x)=c$ pro každé $x\in D$, kde $c$ je pevné reálné číslo. Na $D=\mathbb R$ je zvláštním případem lineární funkce $f(x)=ax+b$ se směrnicí $a=0$ a $b=c$; jejím grafem je vodorovná přímka $y=c$. Na menším oboru jde o příslušnou část tohoto grafu.
\end{definiceBox}
\begin{prikladReseniBox}
Je dáno $f(x)=3$ na $\mathbb R$. Určete $f(-2),f(0),f(5)$, definiční obor, obor hodnot a průsečíky s osami.
\tcblower
Všechny tři hodnoty jsou $3$. $D(f)=\mathbb R$ a $H(f)=\{3\}$. Graf prochází bodem $(0;3)$, což je průsečík s osou $y$. Osu $x$ neprotíná: rovnice $3=0$ nemá řešení. Vodorovná přímka leží tři jednotky nad osou $x$.
\end{prikladReseniBox}
\begin{center}
\begin{tikzpicture}[x=.75cm,y=.7cm,>=stealth]
\draw[step=1,gray!25,very thin] (-4,-1) grid (4,4);
\draw[->] (-4.4,0)--(4.5,0) node[right] {$x$};\draw[->] (0,-1.4)--(0,4.4) node[above] {$y$};
\foreach \x in {-4,-2,2,4}{\draw (\x,.07)--(\x,-.07) node[below] {\small $\x$};}
\foreach \y in {1,2,3,4}{\draw (.07,\y)--(-.07,\y) node[left] {\small $\y$};}
\draw[very thick] (-4,3)--(4,3) node[right] {$y=3$};\fill (0,3) circle (2.5pt);
\end{tikzpicture}
\end{center}
\blok{Vztah k monotonii a k oboru}
Konstantní funkce na intervalu je současně neklesající i nerostoucí: pro $x_1<x_2$ vždy platí $f(x_1)=f(x_2)=c$, tedy obě neostré nerovnosti. Není však přísně rostoucí ani přísně klesající. Pokud $0\in D$, má průsečík s osou $y$ v bodě $(0;c)$; při $c=0$ splývá graf na $\mathbb R$ s osou $x$ a každé $x\in D$ dává funkční hodnotu $0$. Na omezeném oboru kreslíme jen část vodorovné přímky.
\begin{prikladReseniBox}
Funkce $g(x)=-2$ je zadána pro $x\in\langle-1;3)$. Určete její graf, obor hodnot a průsečíky s osami.
\tcblower
Grafem je vodorovná úsečka od plného bodu $(-1;-2)$ k prázdnému bodu $(3;-2)$. $H(g)=\{-2\}$. Protože $0\in\langle-1;3)$, graf protíná osu $y$ v $(0;-2)$. Osu $x$ neprotíná, neboť hodnota je stále $-2$.
\end{prikladReseniBox}
\begin{prikladReseniBox}
Porovnejte rovnice $y=2$ a $x=2$ jako obrázky v souřadnicové rovině. Která z nich udává funkci $y$ proměnné $x$?
\tcblower
$y=2$ je vodorovná přímka: každé reálné $x$ má právě jednu hodnotu $2$, tedy jde o konstantní funkci. $x=2$ je svislá přímka: při argumentu $x=2$ obsahuje nekonečně mnoho různých $y$, zatímco jiná $x$ nemají žádné $y$. Svislicový test selže, takže nejde o graf funkce $y(x)$.
\end{prikladReseniBox}
\poznamka{Nezaměňovat konstantní funkci $y=c$ se svislou přímkou $x=c$. Výrok „konstantní funkce je rostoucí“ je při zde používané přísné definici nepravdivý; správně je neklesající a nerostoucí.}
\blok{Samostatné procvičení}
\begin{zadaniReseniBox}
\begin{enumerate}[label=\alph*)]
\item U $p(x)=0$ na $\mathbb R$ určete obor hodnot a množinu průsečíků s osou $x$.
\item U $q(x)=5$ na $\{-2,0,4\}$ vypište všechny body grafu a obor hodnot.
\item Může mít konstantní funkce $r(x)=7$ při $x=1$ hodnotu $8$? Zdůvodněte.
\end{enumerate}
\tcblower
\begin{enumerate}[label=\alph*)]
\item $H(p)=\{0\}$; celý graf je osa $x$, takže každé $x\in\mathbb R$ dává bod $(x;0)$.
\item Graf tvoří tři body $(-2;5),(0;5),(4;5)$, nikoli souvislá úsečka. $H(q)=\{5\}$.
\item Nemůže, pokud $1\in D(r)$: z předpisu plyne $r(1)=7$. Pokud $1$ v oboru není, hodnota $r(1)$ není zadána vůbec.
\end{enumerate}
\end{zadaniReseniBox}
\begin{zadaniReseniBox}[title={Diagnostika: co skutečně určuje graf?}]
\begin{enumerate}[label=\alph*)]
\item Graf obsahuje body $(-2;4),(0;4),(2;4)$. Dokazují tyto tři body samy, že funkce je konstantní na celé $\mathbb R$?
\item Je předpis $v(x)=5+0x$ na $\mathbb R$ konstantní? Určete jeho směrnici a průsečík s osou $y$.
\end{enumerate}
\tcblower
\begin{enumerate}[label=\alph*)]
\item Nedokazují: mimo uvedené argumenty mohou být hodnoty jiné. Například $w(x)=4+x(x^2-4)$ má v těchto třech bodech hodnotu $4$, ale $w(1)=1$.
\item Ano, $v(x)=5$ pro všechna reálná $x$. Směrnice je $0$ a průsečík s osou $y$ je $(0;5)$.
\end{enumerate}
\end{zadaniReseniBox}
\blok{Volitelná rezerva}
\begin{zadaniReseniBox}
Funkce $s(x)=4$ je dána na neprázdné množině $D=\{-1,1\}$. Je sudá či lichá?
\tcblower
Obor je souměrný vůči nule a $s(-x)=4=s(x)$, proto je sudá. Lichá není: podmínka $s(-x)=-s(x)$ by vyžadovala $4=-4$, což neplatí. Konstantní nulová funkce by na souměrném oboru byla současně sudá i lichá.
\end{zadaniReseniBox}
\blok{Domácí práce}
\begin{zadaniReseniBox}
\begin{enumerate}[label=\alph*)]
\item Pro $f(x)=-3$ na $\mathbb R$ určete $D(f),H(f)$ a průsečíky s osami.
\item Pro $g(x)=2$ na $(-2;1\rangle$ popište konce grafu a vypočtěte $g(0)$.
\item Rozhodněte, zda $y=-1$ a $x=-1$ představují funkci $y(x)$.
\end{enumerate}
\tcblower
\begin{enumerate}[label=\alph*)]
\item $D(f)=\mathbb R$, $H(f)=\{-3\}$. Průsečík s osou $y$ je $(0;-3)$, s osou $x$ žádný.
\item Prázdný bod $(-2;2)$, plný bod $(1;2)$; $g(0)=2$.
\item $y=-1$ je konstantní funkce. $x=-1$ je svislá přímka, která při jednom $x$ připouští mnoho hodnot $y$, takže funkcí $y(x)$ není.
\end{enumerate}
\end{zadaniReseniBox}
\poznamka{Jádro 45 min: motivace a definice 8, graf a úplný rozbor 10, omezený obor 8, srovnání vodorovné a svislé přímky 6, samostatný pokus 7, kontrola a organizace 6 min. Rezerva 5 min, domácí 10–15 min. Další hodina použije lineární funkci jako model reálných veličin.}
\end{document}
